% 解説プリントのコメント（saiten が書く。./tools/kaisetsu.ps1 がこのファイルからプリントを作る）
%% matome
\begin{daiji}
  \item 示量性を調べるときに比べる相手は「同じ系2つ分」$2S'(U,V,N)$ である。差の式の左右を写し間違えると，計算が正しくても何と何を比べたのかが答案から読めなくなる（問題2 (3)）。
  \item 境界条件によらず同じエントロピーになるのは，固定端でも周期境界条件でも，位相空間の面積 $h$ あたりちょうど1つの状態があるからである。(2) で両方の数が $\sqrt{\varepsilon/E_{0}}$ に一致したことが，その1方向版になっている（問題5 (4)，未着手）。
\end{daiji}
\begin{maemo}
  \item 示量性の式の書き間違い：統計力学1-1（旧版）問題4 (3)(4)（2026-10-27，同じ問題）で $S'(2U,2V,2N)-2S(U,V,N)$ とプライムを落とした。今回は (3) で $S'(2U,2V,2N)-S(U,V,N)$ と，プライムと係数2の両方を落とした。同じ日の統計力学1-4 問題2 (4) でも右辺を $\alpha F(T,\alpha V,\alpha N)$ と書いている。示量性の問題では，計算を始める前に「示すこと：$S'(2U,2V,2N)$ と $2S'(U,V,N)$ を比べる」と1行書き，差の式はその行から写す。
  \item 最後の説明の小問の未着手：統計力学1-2 問題3 (5)（2026-11-24）に続いて，今回も問題5 (4) が途中で止まった。止まったら，その小問の前の小問の結果（ここでは (2)）が何を言っているかを言葉にしてみると，手掛かりになることが多い。
\end{maemo}
%% mondai 2
\begin{kotaeran}
  \machigai{3}{○}{記述} 差を $S'(2U,2V,2N)-S(U,V,N)$ と書いた。引くべきものは $2S'(U,V,N)$ で，プライムと係数2が抜けている。次の行の $N\kB[2\log2V-2\log V+\cdots+3-3]$ は $2S'(U,V,N)$ を引いた正しい計算になっていて，答え $2N\kB\log2\neq0$ も正しい。ただし書いた式のとおりに読むと，$N!$ で割った $S$ を $N$ 個分だけ引いたことになり，(4) で $S'$ と $S$ を比べる議論とも混ざる。結論の「$S$ は示量変数でない」も $S'$ のことである。
\end{kotaeran}
\begin{rikai}
  \item 示量性は $S(\lambda U,\lambda V,\lambda N)=\lambda S(U,V,N)$ で，$\lambda=2$ なら「2倍の系のエントロピー」と「元の系2つ分のエントロピー」を比べる。$S'$ の示量性を調べるなら，両辺とも $S'$ である。
  \item 確かめ方：差の式を書いたら，両辺の前の係数がそろっているかを見る。$S'(2U,2V,2N)$ は $2N\kB[\cdots]$，$2S'(U,V,N)$ も $2N\kB[\cdots]$ で，$S(U,V,N)$ だと $N\kB[\cdots]$ となって半分しかない。
\end{rikai}
%% mondai 5
\begin{kotaeran}
  \machigai{4}{×}{未着手} 「固定端では $0<E_{0}n^{2}<\varepsilon$」と書いたところで止まっていて，境界条件によらない理由も，違いが効く場合も書かれていない。
\end{kotaeran}
\begin{rikai}
  \item (2) の結果を運動量で書き直すと (4) の答えになる。$\varepsilon=p^{2}/(2m)$ と $E_{0}=\pi^{2}\hbar^{2}/(2mL^{2})$ から $\sqrt{\varepsilon/E_{0}}=2pL/h$ で，これは「運動量の範囲 $-p$〜$p$ の幅 $2p$」と「位置の範囲 $L$」の積を $h$ で割ったもの，つまり位相空間の面積を $h$ で割ったものである。周期境界条件では運動量が間隔 $h/L$ で正負に並び，固定端の定常波は $\pm p$ の平面波の重ね合わせで $\abs{p}$ が間隔 $h/(2L)$ で正の側だけに並ぶが，1つの定常波を $\pm p$ の両側に広げて数えれば，周期境界条件と同じく位相空間の面積 $h$ あたり1状態になる。
  \item エントロピーは，高いエネルギーでの状態の数（位相空間の体積を $N!\,h^{3N}$ で割ったもの）だけで決まるので，境界条件によらない。違いが効くのは，$\kB T$ が準位の間隔 $E_{0}$ と同程度以下になる場合（小さな容器や極低温）で，このときは基底状態のエネルギー（固定端 $3E_{0}$，周期境界条件 $0$）や縮退度の違いがそのまま熱力学量に出る。
\end{rikai}
